% !PRINTED 266
\weberpage{266}
%<<CONT
section, the interchange $\Bigl(u,\dfrac{\sqrt2}{u}\Bigr)$ carries the
totality of the roots $v$ over into
\[
\Bigl(\frac{2}{p}\Bigr)\frac{\sqrt2}{v}
\]
From this it follows that for $u=0$ all the roots $v$ pass over into
zero, so that none of them becomes infinite; and from this we must
conclude that the $U_1, U_2, \ldots U_{p+1}$ are not only
\gesperrt{polynomials} in $u$, but that each of them must also
contain the factor $u$.

We now conclude first of all, exactly as with the invariant equation
(\S\ 69, 2., 3. and 6.), that
\weq{5}{\Phi_p(v,u)=\Phi_p(u,v),}
and that
\weq{6}{\Phi_p(v,u)=(v^p-u)(v-u^p)+p\sum_{1,p-1}^{h,k}c_{h,k}u^hv^k,}
where $c_{h,k}=c_{k,h}$ are \gesperrt{integers}.
We wish to represent this function in the form
\weq{7}{\Phi_p(v,u)=v^{p+1}+u^{p+1}+\sum_{1,p}^{h,k}a_{h,k}u^hv^k,}
where accordingly $a_{h,k}=a_{k,h}$ are likewise integers, upon whose
divisibility by $p$ nothing further now depends, and in particular
$a_{p,p}=-1$. If in the equation
\[
\Phi_p[f(p\,\omega), f(\omega)]=0
\]
we replace $\omega$ by $\omega+2$, then, on account of the
irreducibility, it results from the relation
\[
f(\omega+2)=e^{-\frac{\pi i}{12}}f(\omega),
\]
that in (7) not \gesperrt{all} terms occur, but only those in which the exponents
$h, k$ satisfy the congruence
\weq{8}{hp+k\equiv p+1\ (\mathrm{mod}\ 24)}

Now we know yet a further property of the functions
$\Phi_p(v,u)$, one which follows from the interchange (9) of the
preceding section already used, where $\varrho=1$ is to be put, and
which, if for brevity we set
\[
\varepsilon=\Bigl(\frac{2}{p}\Bigr)=(-1)^{\frac{p^2-1}{8}}
\]
%CONT>>
